\section{Quasi-coherent sheaves and ample invertible sheaves}
\label{section-ample-quasi-coherent}

\noindent
Theme of this section: in the presence of an ample invertible sheaf
every quasi-coherent sheaf comes from a graded module.

\begin{situation}
\label{situation-ample}
Let $X$ be a scheme.
Let $\mathcal{L}$ be an ample invertible sheaf on $X$.
Set $S = \Gamma_*(X, \mathcal{L})$ as a graded ring.
Set $Y = \text{Proj}(S)$.
Let $f : X \to Y$ be the canonical morphism of Lemma \ref{lemma-map-into-proj}.
It comes equipped with a $\mathbf{Z}$-graded $\mathcal{O}_X$-algebra map
$\bigoplus f^*\mathcal{O}_Y(n) \to \bigoplus \mathcal{L}^{\otimes n}$.
\end{situation}

\noindent
The following lemma is really a special case of the next lemma
but it seems like a good idea to point out its validity first.

\begin{lemma}
\label{lemma-ample-gcd-is-one}
In Situation \ref{situation-ample}.
The canonical morphism $f : X \to Y$
maps $X$ into the open subscheme $W = W_1 \subset Y$
where $\mathcal{O}_Y(1)$ is invertible and where
all multiplication maps
$\mathcal{O}_Y(n) \otimes_{\mathcal{O}_Y} \mathcal{O}_Y(m) \to
\mathcal{O}_Y(n + m)$
are isomorphisms (see
Constructions, Lemma \ref{constructions-lemma-where-invertible}).
Moreover, the maps $f^*\mathcal{O}_Y(n) \to \mathcal{L}^{\otimes n}$
are all isomorphisms.
\end{lemma}

\begin{proof}
By Proposition \ref{proposition-characterize-ample} there exists an integer
$n_0$ such that $\mathcal{L}^{\otimes n}$ is globally generated for all
$n \geq n_0$. Let $x \in X$ be a point. By the above we can find
$a \in S_{n_0}$ and $b \in S_{n_0 + 1}$ such that
$a$ and $b$ do not vanish at $x$. Hence
$f(x) \in D_{+}(a) \cap D_{+}(b) = D_{+}(ab)$. By
Constructions, Lemma \ref{constructions-lemma-where-invertible}
we see that $f(x) \in W_1$ as desired. By
Constructions, Lemma \ref{constructions-lemma-invertible-map-into-proj}
which was used in the construction of the map $f$
the maps
$f^*\mathcal{O}_Y(n_0) \to \mathcal{L}^{\otimes n_0}$ and
$f^*\mathcal{O}_Y(n_0 + 1) \to \mathcal{L}^{\otimes n_0 + 1}$
are isomorphisms in a neighbourhood of $x$. By compatibility with
the algebra structure and the fact that $f$ maps into $W$
we conclude all the maps
$f^*\mathcal{O}_Y(n) \to \mathcal{L}^{\otimes n}$ are isomorphisms
in a neighbourhood of $x$. Hence we win.
\end{proof}

\noindent
Recall from Modules, Definition \ref{modules-definition-gamma-star}
that given a locally ringed space $X$, an invertible sheaf $\mathcal{L}$,
and a $\mathcal{O}_X$-module $\mathcal{F}$ we have the graded
$\Gamma_*(X, \mathcal{L})$-module
$$
\Gamma_*(X, \mathcal{L}, \mathcal{F}) =
\bigoplus\nolimits_{n \in \mathbf{Z}}
\Gamma(X, \mathcal{F} \otimes_{\mathcal{O}_X} \mathcal{L}^{\otimes n}).
$$
The following lemma says that, in Situation \ref{situation-ample},
we can recover a quasi-coherent $\mathcal{O}_X$-module $\mathcal{F}$
from this graded module. Take a look also at
Constructions, Lemma \ref{constructions-lemma-quasi-coherent-projective-space}
where we prove this lemma in the special case $X = \mathbf{P}^n_R$.

\begin{lemma}
\label{lemma-ample-quasi-coherent}
In Situation \ref{situation-ample}.
Let $\mathcal{F}$ be a quasi-coherent sheaf on $X$.
Set $M = \Gamma_*(X, \mathcal{L}, \mathcal{F})$ as a graded $S$-module.
There are isomorphisms
$$
f^*\widetilde{M} \longrightarrow \mathcal{F}
$$
functorial in $\mathcal{F}$ such that
$M_0 \to \Gamma(\text{Proj}(S), \widetilde{M}) \to \Gamma(X, \mathcal{F})$
is the identity map.
\end{lemma}

\begin{proof}
Let $s \in S_{+}$ be homogeneous such that $X_s$ is affine open in $X$.
Recall that $\widetilde{M}|_{D_{+}(s)}$ corresponds to the
$S_{(s)}$-module $M_{(s)}$, see
Constructions, Lemma \ref{constructions-lemma-proj-sheaves}.
Recall that $f^{-1}(D_{+}(s)) = X_s$.
As $X$ carries an ample invertible sheaf it is quasi-compact and
quasi-separated, see Section \ref{section-ample}.
By Lemma \ref{lemma-invert-s-sections} there is a canonical isomorphism
$M_{(s)} = \Gamma_*(X, \mathcal{L}, \mathcal{F})_{(s)} \to
\Gamma(X_s, \mathcal{F})$.
Since $\mathcal{F}$ is quasi-coherent this leads to
a canonical isomorphism
$$
f^*\widetilde{M}|_{X_s} \to \mathcal{F}|_{X_s}
$$
Since $\mathcal{L}$ is ample on $X$ we know that $X$ is covered
by the affine opens of the form $X_s$. Hence it suffices to prove
that the displayed maps glue on overlaps. Proof of this is
omitted.
\end{proof}

\begin{remark}
\label{remark-neurotic}
With assumptions and notation of Lemma \ref{lemma-ample-quasi-coherent}.
Denote the displayed map of the lemma by $\theta_\mathcal{F}$.
Note that the isomorphism $f^*\mathcal{O}_Y(n) \to \mathcal{L}^{\otimes n}$
of Lemma \ref{lemma-ample-gcd-is-one} is just
$\theta_{\mathcal{L}^{\otimes n}}$.
Consider the multiplication maps
$$
\widetilde{M} \otimes_{\mathcal{O}_Y} \mathcal{O}_Y(n)
\longrightarrow
\widetilde{M(n)}
$$
see
Constructions, Equation (\ref{constructions-equation-multiply-more-generally}).
Pull this back to $X$ and consider
$$
\xymatrix{
f^*\widetilde{M} \otimes_{\mathcal{O}_X} f^*\mathcal{O}_Y(n)
\ar[r]
\ar[d]_{\theta_\mathcal{F} \otimes \theta_{\mathcal{L}^{\otimes n}}}
&
f^*\widetilde{M(n)}
\ar[d]^{\theta_{\mathcal{F} \otimes \mathcal{L}^{\otimes n}}}
\\
\mathcal{F} \otimes \mathcal{L}^{\otimes n} \ar[r]^{\text{id}} &
\mathcal{F} \otimes \mathcal{L}^{\otimes n}
}
$$
Here we have used the obvious identification
$M(n) = \Gamma_*(X, \mathcal{L}, \mathcal{F} \otimes \mathcal{L}^{\otimes n})$.
This diagram commutes. Proof omitted.
\end{remark}

\noindent
It should be possible to deduce the following lemma from
Lemma \ref{lemma-ample-quasi-coherent} (or conversely) but it seems
simpler to just repeat the proof.

\begin{lemma}
\label{lemma-proj-quasi-coherent}
Let $S$ be a graded ring such that $X = \text{Proj}(S)$ is quasi-compact.
Let $\mathcal{F}$ be a quasi-coherent $\mathcal{O}_X$-module. Set
$M = \bigoplus_{n \in \mathbf{Z}} \Gamma(X, \mathcal{F}(n))$ as
a graded $S$-module, see
Constructions, Section \ref{constructions-section-invertible-on-proj}.
The map
$$
\widetilde{M} \longrightarrow \mathcal{F}
$$
of Constructions, Lemma
\ref{constructions-lemma-comparison-proj-quasi-coherent}
is an isomorphism.
If $X$ is covered by standard opens $D_+(f)$ where $f$ has degree $1$,
then the induced maps
$M_n \to \Gamma(X, \mathcal{F}(n))$ are the identity maps.
\end{lemma}

\begin{proof}
Since $X$ is quasi-compact we can find homogeneous elements
$f_1, \ldots, f_n \in S$ of positive degrees such that
$X = D_+(f_1) \cup \ldots \cup D_+(f_n)$. Let $d$ be the
least common multiple of the degrees of $f_1, \ldots, f_n$.
After replacing $f_i$ by a power we may assume that each
$f_i$ has degree $d$. Then we see that $\mathcal{L} = \mathcal{O}_X(d)$ is
invertible, the multiplication maps
$\mathcal{O}_X(ad) \otimes \mathcal{O}_X(bd) \to \mathcal{O}_X((a + b)d)$
are isomorphisms, and each $f_i$ determines a global section $s_i$
of $\mathcal{L}$ such that $X_{s_i} = D_+(f_i)$, see
Constructions, Lemmas \ref{constructions-lemma-where-invertible} and
\ref{constructions-lemma-principal-open}.
Thus $\Gamma(X, \mathcal{F}(ad)) =
\Gamma(X, \mathcal{F} \otimes \mathcal{L}^{\otimes a})$.
Recall that $\widetilde{M}|_{D_{+}(f_i)}$ corresponds to the
$S_{(f_i)}$-module $M_{(f_i)}$, see
Constructions, Lemma \ref{constructions-lemma-proj-sheaves}.
Since the degree of $f_i$ is $d$, the isomorphism class of
$M_{(f_i)}$ depends only on the homogeneous summands of $M$ of
degree divisible by $d$. More precisely, the isomorphism class of
$M_{(f_i)}$ depends only on the graded $\Gamma_*(X, \mathcal{L})$-module
$\Gamma_*(X, \mathcal{L}, \mathcal{F})$
and the image $s_i$ of $f_i$ in $\Gamma_*(X, \mathcal{L})$.
The scheme $X$ is quasi-compact by assumption and
separated by Constructions, Lemma \ref{constructions-lemma-proj-separated}.
By Lemma \ref{lemma-invert-s-sections} there is a canonical isomorphism
$$
M_{(f_i)} = \Gamma_*(X, \mathcal{L}, \mathcal{F})_{(s_i)} \to
\Gamma(X_{s_i}, \mathcal{F}).
$$
The construction of the map in Constructions, Lemma
\ref{constructions-lemma-comparison-proj-quasi-coherent}
then shows that it is an isomorphism over $D_+(f_i)$
hence an isomorphism as $X$ is covered by these opens.
We omit the proof of the final statement.
\end{proof}




